% GENERATED FILE. Edit canonical lessons, then run npm run generate:books.
% canonical-source-hash: fnv1a64:75003d087dfce844
% canonical-lessons: UR-C409-janib, UR-C409-rukh, UR-C409-nazdik, UR-C409-gosha, UR-C409-jhankna, UR-R409-review-402-409

\chapter{Which Way?}
\label{ch:ur-which-way}
\hlchaptermodality{409}

\begin{chapteropening}
\textbf{By the end of this chapter:} \emph{I can say which way something faces and what lies close by.}

Complete the last lesson of chapter 409: I can say which way something faces and what lies close by.
\end{chapteropening}

\section[\texorpdfstring{\ur{جانب} (jānib)}{جانب (jānib)}]{\ur{جانب} (jānib) — towards; a side (formal)}

\label{lesson:UR-C409-janib}



\begin{quote}
Before the new one: say the Urdu for the east, then the Urdu for the west.
\end{quote}

\subsection*{You'll want to know: \ur{جانب}}
\textbf{\ur{جانب}} — \emph{jānib} — \textquotedblleft{}towards; a side (formal)\textquotedblright{}.

\begin{quote}
\ur{دریا} \ur{کی} \ur{جانب} \ur{ایک} \ur{راستہ} \ur{جاتا} \ur{ہے۔}
\end{quote}

\begin{quote}
\emph{daryā kī jānib ek rāstā jātā hai.}
\end{quote}

\begin{quote}
A path leads towards the river.
\end{quote}

Feminine.

\begin{cousinweb}
\textbf{\ur{جانب}} (\emph{jānib}) is usually traced, through Persian, to Arabic \emph{jānib}, \textquotedblleft{}a side\textquotedblright{}; \textbf{\ur{کی} \ur{جانب}} means \textquotedblleft{}towards\textquotedblright{}, a more formal \textbf{\ur{کی} \ur{طرف}}.
\end{cousinweb}

\begin{grammarlens}[title={What you've built}]
One more word for getting around.
\end{grammarlens}

\subsection*{Your turn}
\begin{itemize}
\raggedright
  \item \emph{Say it:} \emph{jānib}
  \item \emph{Say it:} \emph{jānib}, once more
  \item \emph{Say it:} \emph{daryā kī jānib ek rāstā jātā hai.} — A path leads towards the river.
  \item \emph{Recall:} say the Urdu for the east, then the Urdu for the west, then say \emph{jānib} again
\end{itemize}

\begin{tcolorbox}[breakable,colback=teal!4,colframe=teal!35!black,title={Before you move on}]
What does \ur{جانب} mean? (\textquotedblleft{}Towards; a side (formal)\textquotedblright{}.) Say it once more.
\end{tcolorbox}
\section[\texorpdfstring{\ur{رخ} (rukh)}{رخ (rukh)}]{\ur{رخ} (rukh) — the way something faces; a direction}

\label{lesson:UR-C409-rukh}



\begin{quote}
Before the new one: say the Urdu for the west, then the Urdu for towards.
\end{quote}

\subsection*{You'll want to know: \ur{رخ}}
\textbf{\ur{رخ}} — \emph{rukh} — \textquotedblleft{}the way something faces; a direction\textquotedblright{}.

\begin{quote}
\ur{ہوا} \ur{کا} \ur{رخ} \ur{شمال} \ur{کی} \ur{طرف} \ur{ہے۔}
\end{quote}

\begin{quote}
\emph{hawā kā rukh shimāl kī taraf hai.}
\end{quote}

\begin{quote}
The wind is blowing towards the north.
\end{quote}

Masculine.

\begin{cousinweb}
is usually traced to Persian \emph{rukh}, \textquotedblleft{}the face, a cheek\textquotedblright{}: the way a thing faces is its \textbf{\ur{رخ}}.
\end{cousinweb}

\begin{grammarlens}[title={What you've built}]
One more word for getting around.
\end{grammarlens}

\subsection*{Your turn}
\begin{itemize}
\raggedright
  \item \emph{Say it:} \emph{rukh}
  \item \emph{Say it:} \emph{rukh}, once more
  \item \emph{Say it:} \emph{hawā kā rukh shimāl kī taraf hai.} — The wind is blowing towards the north.
  \item \emph{Recall:} say the Urdu for the west, then the Urdu for towards, then say \emph{rukh} again
\end{itemize}

\begin{tcolorbox}[breakable,colback=teal!4,colframe=teal!35!black,title={Before you move on}]
What does \ur{رخ} mean? (\textquotedblleft{}The way something faces; a direction\textquotedblright{}.) Say it once more.
\end{tcolorbox}
\section[\texorpdfstring{\ur{نزدیک} (nazdīk)}{نزدیک (nazdīk)}]{\ur{نزدیک} (nazdīk) — close by; near}

\label{lesson:UR-C409-nazdik}



\begin{quote}
Before the new one: say the Urdu for towards, then the Urdu for the way something faces.
\end{quote}

\subsection*{You'll want to know: \ur{نزدیک}}
\textbf{\ur{نزدیک}} — \emph{nazdīk} — \textquotedblleft{}close by; near\textquotedblright{}.

\begin{quote}
\ur{سرائے} \ur{اسٹیشن} \ur{کے} \ur{نزدیک} \ur{ہے۔}
\end{quote}

\begin{quote}
\emph{sarāe isṭeshan ke nazdīk hai.}
\end{quote}

\begin{quote}
The inn is close to the station.
\end{quote}

\begin{cousinweb}
\textbf{\ur{نزدیک}} (\emph{nazdīk}) is usually traced to Persian \emph{nazdīk}, \textquotedblleft{}near\textquotedblright{}; \textbf{\ur{قریب}} says the same.
\end{cousinweb}

\begin{grammarlens}[title={What you've built}]
One more word for getting around.
\end{grammarlens}

\subsection*{Your turn}
\begin{itemize}
\raggedright
  \item \emph{Say it:} \emph{nazdīk}
  \item \emph{Say it:} \emph{nazdīk}, once more
  \item \emph{Say it:} \emph{sarāe isṭeshan ke nazdīk hai.} — The inn is close to the station.
  \item \emph{Recall:} say the Urdu for towards, then the Urdu for the way something faces, then say \emph{nazdīk} again
\end{itemize}

\begin{tcolorbox}[breakable,colback=teal!4,colframe=teal!35!black,title={Before you move on}]
What does \ur{نزدیک} mean? (\textquotedblleft{}Close by; near\textquotedblright{}.) Say it once more.
\end{tcolorbox}
\section[\texorpdfstring{\ur{گوشہ} (goshā)}{گوشہ (goshā)}]{\ur{گوشہ} (goshā) — a nook; a corner}

\label{lesson:UR-C409-gosha}



\begin{quote}
Before the new one: say the Urdu for the way something faces, then the Urdu for close by.
\end{quote}

\subsection*{You'll want to know: \ur{گوشہ}}
\textbf{\ur{گوشہ}} — \emph{goshā} — \textquotedblleft{}a nook; a corner\textquotedblright{}.

\begin{quote}
\ur{کمرے} \ur{کے} \ur{ایک} \ur{گوشے} \ur{میں} \ur{میز} \ur{ہے۔}
\end{quote}

\begin{quote}
\emph{kamre ke ek goshe meṅ mez hai.}
\end{quote}

\begin{quote}
There is a table in one corner of the room.
\end{quote}

Masculine.

\begin{cousinweb}
\textbf{\ur{گوشہ}} (\emph{goshā}) is usually traced to Persian \emph{gosha}, \textquotedblleft{}a corner\textquotedblright{}; it is a more literary word than \textbf{\ur{کونا}}.
\end{cousinweb}

\begin{grammarlens}[title={What you've built}]
One more word for getting around.
\end{grammarlens}

\subsection*{Your turn}
\begin{itemize}
\raggedright
  \item \emph{Say it:} \emph{goshā}
  \item \emph{Say it:} \emph{goshā}, once more
  \item \emph{Say it:} \emph{kamre ke ek goshe meṅ mez hai.} — There is a table in one corner of the room.
  \item \emph{Recall:} say the Urdu for the way something faces, then the Urdu for close by, then say \emph{goshā} again
\end{itemize}

\begin{tcolorbox}[breakable,colback=teal!4,colframe=teal!35!black,title={Before you move on}]
What does \ur{گوشہ} mean? (\textquotedblleft{}A nook; a corner\textquotedblright{}.) Say it once more.
\end{tcolorbox}
\section[\texorpdfstring{\ur{جھانکنا} (jhāṅknā)}{جھانکنا (jhāṅknā)}]{\ur{جھانکنا} (jhāṅknā) — to peep; to peer (in)}

\label{lesson:UR-C409-jhankna}



\begin{quote}
Before the new one: say the Urdu for close by, then the Urdu for a nook.
\end{quote}

\subsection*{You'll want to know: \ur{جھانکنا}}
\textbf{\ur{جھانکنا}} — \emph{jhāṅknā} — \textquotedblleft{}to peep; to peer (in)\textquotedblright{}.

\begin{quote}
\ur{سارا} \ur{نے} \ur{دروازے} \ur{سے} \ur{اندر} \ur{جھانکا۔}
\end{quote}

\begin{quote}
\emph{sārā ne darvāze se andar jhāṅkā.}
\end{quote}

\begin{quote}
Sara peeped in round the door.
\end{quote}

Past, with \ur{نے} (he or she …): \ur{اس} \ur{نے} \ur{کھڑکی} \ur{سے} \ur{جھانکا}, \emph{us ne khiṛkī se jhāṅkā}.

\begin{cousinweb}
\textbf{\ur{جھانکنا}} (\emph{jhāṅknā}) is inherited: CDIAL lists the Hindi cousin \emph{jhā̃khnā}, \textquotedblleft{}to peep at\textquotedblright{}, under \emph{jhaṅkh}, a form it reconstructs (CDIAL 5325).
\end{cousinweb}

\begin{grammarlens}[title={What you've built}]
One more word for getting around.
\end{grammarlens}

\subsection*{Your turn}
\begin{itemize}
\raggedright
  \item \emph{Say it:} \emph{jhāṅknā}
  \item \emph{Say it:} \emph{jhāṅknā}, once more
  \item \emph{Say it:} \emph{sārā ne darvāze se andar jhāṅkā.} — Sara peeped in round the door.
  \item \emph{Recall:} say the Urdu for close by, then the Urdu for a nook, then say \emph{jhāṅknā} again
\end{itemize}

\begin{tcolorbox}[breakable,colback=teal!4,colframe=teal!35!black,title={Before you move on}]
What does \ur{جھانکنا} mean? (\textquotedblleft{}To peep; to peer (in)\textquotedblright{}.) Say it once more.
\end{tcolorbox}
\section[(dialogue)]{Looking back: from Long Journeys to Which Way?}

\label{lesson:UR-R409-review-402-409}



\begin{quote}
Say the Urdu for a distance travelled and a guide.
\end{quote}

\subsection*{Your turn}
Say these aloud:

\begin{itemize}
\raggedright
  \item a distance travelled, a guide, to tire (someone) out, to put (someone) to sleep, to wake (someone) up
  \item an inn, a stay, a place to stay, to spread out (a sheet), a lift
  \item a washroom, soap, a water heater, to drip, a facility
  \item a register (book), identification, a doorman, to send for, a tip (for service)
  \item wood, to kindle (a fire), to wrap oneself in, a canopy, a night's shelter
  \item out of order, to get broken, noise, a repair, to move (something) aside
  \item a direction, the north, the south, the east, the west
  \item towards, the way something faces, close by, a nook, to peep
\end{itemize}

\begin{tcolorbox}[breakable,colback=teal!4,colframe=teal!35!black,title={Before you move on}]
A distance travelled? (\textbf{\ur{مسافت}}, \emph{masāfat}.) A guide? (\textbf{\ur{رہنما}}, \emph{rahnumā}.) To tire (someone) out? (\textbf{\ur{تھکانا}}, \emph{thakānā}.) To put (someone) to sleep? (\textbf{\ur{سلانا}}, \emph{sulānā}.) To wake (someone) up? (\textbf{\ur{جگانا}}, \emph{jagānā}.) An inn? (\textbf{\ur{سرائے}}, \emph{sarāe}.) A stay? (\textbf{\ur{قیام}}, \emph{qiyām}.) A place to stay? (\textbf{\ur{ٹھکانا}}, \emph{ṭhikānā}.) To spread out (a sheet)? (\textbf{\ur{بچھانا}}, \emph{bichhānā}.) A lift? (\textbf{\ur{لفٹ}}, \emph{lifṭ}.) A washroom? (\textbf{\ur{غسل} \ur{خانہ}}, \emph{ghusl khānā}.) Soap? (\textbf{\ur{صابن}}, \emph{sābun}.) A water heater? (\textbf{\ur{گیزر}}, \emph{gīzar}.) To drip? (\textbf{\ur{ٹپکنا}}, \emph{ṭapaknā}.) A facility? (\textbf{\ur{سہولت}}, \emph{sahūlat}.) A register (book)? (\textbf{\ur{رجسٹر}}, \emph{rajisṭar}.) Identification? (\textbf{\ur{شناخت}}, \emph{shanākht}.) A doorman? (\textbf{\ur{دربان}}, \emph{darbān}.) To send for? (\textbf{\ur{منگانا}}, \emph{mangānā}.) A tip (for service)? (\textbf{\ur{بخشش}}, \emph{bakhshish}.) Wood? (\textbf{\ur{لکڑی}}, \emph{lakṛī}.) To kindle (a fire)? (\textbf{\ur{سلگانا}}, \emph{sulgānā}.) To wrap oneself in? (\textbf{\ur{اوڑھنا}}, \emph{oṛhnā}.) A canopy? (\textbf{\ur{شامیانہ}}, \emph{shāmiyānā}.) A night's shelter? (\textbf{\ur{بسیرا}}, \emph{baserā}.) Out of order? (\textbf{\ur{خراب}}, \emph{kharāb}.) To get broken? (\textbf{\ur{ٹوٹنا}}, \emph{ṭūṭnā}.) Noise? (\textbf{\ur{شور}}, \emph{shor}.) A repair? (\textbf{\ur{مرمت}}, \emph{marammat}.) To move (something) aside? (\textbf{\ur{ہٹانا}}, \emph{haṭānā}.) A direction? (\textbf{\ur{سمت}}, \emph{simt}.) The north? (\textbf{\ur{شمال}}, \emph{shimāl}.) The south? (\textbf{\ur{جنوب}}, \emph{junūb}.) The east? (\textbf{\ur{مشرق}}, \emph{mashriq}.) The west? (\textbf{\ur{مغرب}}, \emph{maghrib}.) Towards? (\textbf{\ur{جانب}}, \emph{jānib}.) The way something faces? (\textbf{\ur{رخ}}, \emph{rukh}.) Close by? (\textbf{\ur{نزدیک}}, \emph{nazdīk}.) A nook? (\textbf{\ur{گوشہ}}, \emph{goshā}.) To peep? (\textbf{\ur{جھانکنا}}, \emph{jhāṅknā}.)
\end{tcolorbox}
